\section{Introduction}

Two pictures of an agent network are in circulation. In the first, agents discover
each other in the open: they publish what they can do, broadcast what they need, and
are matched by relevance rather than by prior acquaintance. In the second, agents
act for people who already work together: membership is bounded, interaction repeats,
context accumulates, and someone is accountable when the work is wrong. The first
picture is usually called public and the second private, and the difference is
usually described as a difference in \emph{visibility}.

Visibility is the wrong frame, and a deployed system makes that obvious. Whether a
stranger can \emph{find} an agent and whether that stranger can \emph{read what the
agent holds} are separately enforced, and they do not co-vary. An open directory can
front a counterpart that answers questions and reveals nothing; a bounded roster can
front a counterpart whose entire working store is readable. The interesting question
is not who can be seen but what crosses the boundary once contact is made---and,
prior to that, whether the relational language the field uses corresponds to anything
the runtime enforces at all.

That last point is the paper's subject. Agent networks encode relationships in two
places. Some of it is \arch{architecture}: the directory a requester may enumerate,
whether a counterpart exposes an interface or a store, whether state persists between
encounters, whether identity is stable across them. A kernel checks these; an agent
cannot talk its way past them. The rest is \txt{text}: a trust score on a card, a
stated shared objective, a note that two parties have collaborated before, a claim
that answers are logged and attributed. These reach the model as tokens, and whether
they change behavior is an empirical question that, to our knowledge, has not been
asked.

We suspect the answer is uncomfortable. Preliminary work in our own setting suggests
that once reach is equalized between a public and a private configuration, most of
the measured difference between them disappears---and reach is the crudest of the
architectural properties. If the textual group is largely inert on top of that, then
much of what agent-network design currently invests in---richer cards, reputation
fields, declared alignment, relationship history surfaced as context---is decoration
that a runtime does not act on.

\paragraph{Contributions.}
\begin{enumerate}[leftmargin=1.4em,itemsep=1pt,topsep=2pt]
\item A definition of public and private agent networks over nine relational
      dimensions, partitioned into an \arch{architectural} group a kernel enforces
      and a \txt{textual} group that exists only as seeded context
      (Section~\ref{sec:framework}).
\item An operationalization of the architectural group as a $2\times2$ factorial over
      \emph{discovery} and \emph{access}, run as configurations of a single
      capability kernel so that a condition differs from another in named values
      rather than in implementation (Section~\ref{sec:operational}).
\item A seeding protocol for the textual group, with an explicit validity boundary:
      we measure whether agents act on asserted relational context, not what trust
      does in a human network (Section~\ref{sec:seeding}).
\item A measurement design whose primary metrics require no judge---deliverables are
      hidden-profile artifacts with forced values---and whose three judged secondary
      metrics are blind, pairwise, cross-family, and reported with agreement
      (Section~\ref{sec:measurement}).
\item Seven predictions chosen for being ones a careful reader would not confidently
      bet on in advance, together with the falsification condition for each
      (Section~\ref{sec:design}).
\end{enumerate}

\paragraph{What we deliberately do not claim.}
We do not study how a public stranger becomes a private teammate; the transition is
instrumented in the deployed system this design is drawn from, but a within-pair
before/after study is a separate paper. We make no safety claim about either regime:
both are deny-by-default for action, and we report leakage as a measurement rather
than a verdict. And we state in advance that the \emph{existence} of a
task-dependent crossing between regimes is a property of how we allocate facts, not
a discovery; what is measured is where it falls, what it costs, and how much of the
attainable each configuration realizes.
